% =====================================================================
%  Memoria de la práctica de negociación automática (NB3)
%  Sistemas Multiagente — UIMP
%
%  Compilar (dos pasadas para referencias e índice):
%      pdflatex memoria.tex && pdflatex memoria.tex
%  o subir la carpeta docs/memoria/ a Overleaf.
% =====================================================================
\documentclass[11pt,a4paper]{article}

% --- Idioma y codificación -------------------------------------------
\usepackage[utf8]{inputenc}
\usepackage[T1]{fontenc}
\usepackage{lmodern}
% english se carga para citar el enunciado original; spanish es el idioma principal.
\usepackage[english,spanish,es-noquoting,es-tabla,es-nodecimaldot]{babel}
\usepackage{csquotes}

% --- Maquetación -----------------------------------------------------
\usepackage[a4paper,margin=2.5cm]{geometry}
\usepackage{microtype}
\usepackage{parskip}
\usepackage{enumitem}
\setlist{itemsep=0.3em}

% --- Tablas ----------------------------------------------------------
\usepackage{booktabs}
\usepackage{tabularx}
\usepackage{array}
\newcolumntype{Y}{>{\raggedright\arraybackslash}X}
\renewcommand{\arraystretch}{1.15}

% --- Matemáticas -----------------------------------------------------
\usepackage{amsmath}
\usepackage{amssymb}

% --- Código y logs ---------------------------------------------------
\usepackage{xcolor}
\usepackage{listings}
\definecolor{codebg}{gray}{0.95}
\lstset{
  basicstyle=\ttfamily\footnotesize,
  backgroundcolor=\color{codebg},
  frame=single,
  rulecolor=\color{gray!50},
  breaklines=true,
  breakatwhitespace=false,
  columns=fullflexible,
  keepspaces=true,
  showstringspaces=false,
  literate={×}{{$\times$}}1 {á}{{\'a}}1 {é}{{\'e}}1 {í}{{\'i}}1 {ó}{{\'o}}1 {ú}{{\'u}}1 {ñ}{{\~n}}1
}

% --- Enlaces ---------------------------------------------------------
\usepackage[hidelinks]{hyperref}
\usepackage{url}

% --- Macros propias --------------------------------------------------
% Texto de código en línea (nombres de clases, ficheros, mensajes...).
\newcommand{\code}[1]{\texttt{#1}}
% Enlace a una ruta del repositorio en GitHub.
\newcommand{\repourl}{https://github.com/migueldlfuentem/nb-agent/tree/main}
\newcommand{\repo}[2]{\href{\repourl/#1}{#2}}

% --- Portada ---------------------------------------------------------
\title{%
  \textbf{Práctica de Negociación Automática}\\[0.3em]
  \large Agente negociador basado en el algoritmo NB\textsuperscript{3}\\[0.2em]
  \large Sistemas Multiagente --- UIMP}
\author{Nombre Apellidos}  % TODO: completar
\date{\today}

\begin{document}

\maketitle

\begin{abstract}
  Esta memoria recoge la resolución de los ejercicios de la sección~5 del manual
  \emph{Implementing an NB\textsuperscript{3} Agent} (D.~de~Jonge, 2015), sobre el dominio
  \emph{Commodity Market}: cinco agentes que intercambian cinco materias primas mediante el
  protocolo de negociación no estructurado. El código fuente, los logs de las sesiones y esta
  memoria están disponibles en el repositorio
  \href{https://github.com/migueldlfuentem/nb-agent}{\nolinkurl{github.com/migueldlfuentem/nb-agent}}.
\end{abstract}

\tableofcontents
\newpage

% =====================================================================
%  Ejercicio 1: análisis de las sesiones del agente NB3 de ejemplo
%  (versión LaTeX de docs/ejercicio1/README.md)
% =====================================================================
\section{\texorpdfstring{Ejercicio 1: análisis de las sesiones de negociación del agente NB\textsuperscript{3} de ejemplo}{Ejercicio 1: análisis de las sesiones de negociación del agente NB3 de ejemplo}}
\label{sec:ej1}

% ---------------------------------------------------------------------
\subsection*{Enunciado}
\addcontentsline{toc}{subsection}{Enunciado}

\begin{otherlanguage}{english}
\begin{quote}
  \textit{Manual NB\textsuperscript{3}, section 5: Exercises, Exercise 1}

  Start a negotiation session by running the class \code{RunMarketAndFiveAgents}. After 30 seconds
  you should see the results of the negotiations.

  Now go to the \code{logfolderPath} specified in the source code of this class and check that
  indeed a folder has been created with the log files for all negotiators and the server.

  Open the server log file and determine how many proposals each agent has made. For each proposal
  message sent, the log file should contain something like:

\begin{lstlisting}
sent message:
message ID: Bob 1
conversation ID:
from: Bob
to: [Charles]
particle: PROPOSE
content: Proposal ID: Bob1
Participants: [Bob, Charles]
Proposed Deal: [agent 2 gives 7 units of 1 to agent 3, agent 3 gives 10 units of 0 to agent 2, agent 3 gives 10 units of 2 to agent 2]
\end{lstlisting}

  Of course, the details are different for each proposal. The line \code{particle: PROPOSE}
  indicates that this is a proposal. The log will also contain messages of type \code{REGISTER},
  \code{INFORM} and \code{START}. You can ignore these, they form part of the communication protocol.

  If the negotiations were successful you will also find messages of type \code{ACCEPT} and
  \code{CONFIRM}. \textbf{How many times has a proposal been accepted? How many proposals have been
  confirmed?}

  Remember that the instance of the CommodityMarket is created randomly every time you run a new
  negotiation session. Therefore, in some cases the agents may not find any good proposal at all, so
  you may need to run several negotiation sessions before anything interesting happens. If after a
  couple of times the agents still don't make any proposals it may be better to continue with the
  other exercises to improve the agent and get back to this exercise with your improved agent.
\end{quote}
\end{otherlanguage}

Por coherencia con la práctica, a lo largo de esta documentación se mantienen en inglés los nombres
de las materias primas (\emph{Gold}, \emph{Oil}, \emph{Iron}, \emph{Grain} y \emph{Wood}).

% ---------------------------------------------------------------------
\subsection*{Sesiones realizadas}
\addcontentsline{toc}{subsection}{Sesiones realizadas}

Se han realizado \textbf{7 sesiones} independientes de 30\,s. Cada una genera una instancia
aleatoria distinta del mercado (activos y requisitos). En cada sesión se crea la carpeta de logs con
un \code{server.log}, y dos logs por agente: \code{<Agente>\_nego.log} y \code{<Agente>\_search.log}.

Los logs de cada sesión están disponibles en el repositorio, en
\repo{docs/ejercicio1/logs}{\code{docs/ejercicio1/logs/}}:

\begin{itemize}
  \item \code{server.log}: todos los mensajes intercambiados, los activos y requisitos iniciales y
        las decisiones del Notario.
  \item \code{console.log}: salida de la sesión, que incluye la tabla final de utilidades
        (antes/después).
  \item \code{<Agente>\_nego.log}: decisiones de cada agente (\code{WAIT}, \code{PROPOSE},
        \code{ACCEPT}) y sus niveles de aspiración.
\end{itemize}

Los \code{<Agente>\_search.log} (árbol de búsqueda) no se incluyen porque ocupan entre 20 y 80\,MB
por agente y sesión.

% ---------------------------------------------------------------------
\subsection{Cómo se han contado los mensajes}
\label{sec:ej1-conteo}

El enunciado sugiere contar las propuestas buscando en el log las entradas con
\code{particle: PROPOSE}. Sin embargo, al analizar los logs se ha detectado que \textbf{sumar las
ocurrencias de \code{PROPOSE} no es válido, porque duplica registros}. El servidor escribe cada
mensaje varias veces (\emph{broadcast}): una al recibirlo del emisor (\code{received message:}) y
otra por cada destinatario al que lo reenvía (\code{sent message:}). Una misma propuesta aparece,
por tanto, como mínimo dos veces, y más si va dirigida a varios agentes. La diferencia es
considerable (tabla~\ref{tab:ej1-conteo}).

\begin{table}[htbp]
  \centering
  \caption{Propuestas reales frente a las entradas que aparecen en \code{server.log}.}
  \label{tab:ej1-conteo}
  \begin{tabular}{cccc}
    \toprule
    Sesión & Propuestas reales            & Entradas            & Total de líneas         \\
           & (\code{received message})    & \code{sent message} & \code{particle: PROPOSE} \\
    \midrule
    1 & 2  & 2  & 4  \\
    4 & 38 & 49 & 87 \\
    6 & 31 & 33 & 64 \\
    \bottomrule
  \end{tabular}
\end{table}

Ni siquiera contar solo los bloques \code{sent message}, como en el ejemplo del enunciado, da el
número correcto. En la sesión~4 había propuestas a tres bandas (\code{[Bob, Eve, Charles]}) que se
reenviaron a dos destinatarios y figuran dos veces. Por ello se ha seguido este criterio:

\begin{itemize}
  \item \textbf{PROPOSE} y \textbf{ACCEPT} se cuentan solo en los bloques \code{received message:}:
        cada bloque corresponde a un único mensaje enviado por un agente.
  \item \textbf{CONFIRM} e \textbf{ILLEGAL} se cuentan como \textbf{IDs de propuesta distintos},
        porque los envía el \code{NOTARY} una vez por destinatario.
  \item \code{agent N} se traduce al nombre del agente con el orden de la tabla
        \code{REQUIREMENTS}. Ese orden es el del registro en el Notario, y sirve también como mapa
        ID~$\rightarrow$~nombre: es el array \code{agentNames} con el que se inicializa el algoritmo
        NB\textsuperscript{3}.
\end{itemize}

% ---------------------------------------------------------------------
\subsection{Resultados}
\label{sec:ej1-resultados}

Esta sección presenta los datos extraídos de los logs de las 7 sesiones: la actividad de cada
sesión, las propuestas de cada agente y los acuerdos alcanzados. Aquí se describen los resultados;
sus causas se analizan en la sección~\ref{sec:ej1-interpretacion}.

\subsubsection{Resumen por sesión}
\label{sec:ej1-resumen}

La tabla~\ref{tab:ej1-resumen} recoge, para cada sesión, el número de propuestas realizadas, los
mensajes de aceptación (\code{ACCEPT}), los acuerdos confirmados por el Notario (\code{CONFIRM}) y las
propuestas declaradas ilegales (\code{ILLEGAL}), junto con un enlace a sus logs.

\begin{table}[htbp]
  \centering
  \caption{Resumen de las 7 sesiones de negociación.}
  \label{tab:ej1-resumen}
  \begin{tabular}{clcccc}
    \toprule
    Sesión & Logs & Propuestas & ACCEPT & Confirmados & ILLEGAL \\
    \midrule
    1 & \href{https://github.com/migueldlfuentem/nb-agent/tree/main/docs/ejercicio1/logs/sesion1_2026-11-14__20-28-22}{\code{sesion1}} & 2  & 1  & 1 & 0 \\
    2 & \href{https://github.com/migueldlfuentem/nb-agent/tree/main/docs/ejercicio1/logs/sesion2_2026-11-14__20-28-56}{\code{sesion2}} & 0  & 0  & 0 & 0 \\
    3 & \href{https://github.com/migueldlfuentem/nb-agent/tree/main/docs/ejercicio1/logs/sesion3_2026-11-14__20-29-31}{\code{sesion3}} & 0  & 0  & 0 & 0 \\
    4 & \href{https://github.com/migueldlfuentem/nb-agent/tree/main/docs/ejercicio1/logs/sesion4_2026-11-14__20-30-05}{\code{sesion4}} & 38 & 12 & 2 & 0 \\
    5 & \href{https://github.com/migueldlfuentem/nb-agent/tree/main/docs/ejercicio1/logs/sesion5_2026-11-14__20-30-41}{\code{sesion5}} & 15 & 1  & 1 & 0 \\
    6 & \href{https://github.com/migueldlfuentem/nb-agent/tree/main/docs/ejercicio1/logs/sesion6_2026-11-14__20-31-14}{\code{sesion6}} & 31 & 3  & 1 & 1 \\
    7 & \href{https://github.com/migueldlfuentem/nb-agent/tree/main/docs/ejercicio1/logs/sesion7_2026-11-14__20-32-49}{\code{sesion7}} & 0  & 0  & 0 & 0 \\
    \midrule
    \textbf{Total} & & \textbf{86} & \textbf{17} & \textbf{5} & \textbf{1} \\
    \bottomrule
  \end{tabular}
\end{table}

La actividad varía mucho de una sesión a otra: en 3 de las 7 no se hizo ninguna propuesta, mientras
que las sesiones 4 y 6 concentran 69 de las 86. Además, el número de acuerdos confirmados es muy
inferior al de aceptaciones (5 frente a 17), y una propuesta llegó a declararse ilegal (ver
secciones~\ref{sec:ej1-sin-actividad} y~\ref{sec:ej1-aceptaciones}).

\subsubsection{Propuestas por agente}
\label{sec:ej1-por-agente}

La tabla~\ref{tab:ej1-por-agente} desglosa las propuestas de la tabla anterior según el agente que
las realizó, con una columna por sesión (S1--S7) y el total acumulado de cada agente.

\begin{table}[htbp]
  \centering
  \caption{Propuestas realizadas por cada agente en cada sesión.}
  \label{tab:ej1-por-agente}
  \begin{tabular}{lcccccccc}
    \toprule
    Agente & S1 & S2 & S3 & S4 & S5 & S6 & S7 & \textbf{Total} \\
    \midrule
    Alice   & 0 & 0 & 0 & 0  & 15 & 0  & 0 & \textbf{15} \\
    Bob     & 0 & 0 & 0 & 23 & 0  & 0  & 0 & \textbf{23} \\
    Charles & 1 & 0 & 0 & 0  & 0  & 24 & 0 & \textbf{25} \\
    David   & 1 & 0 & 0 & 14 & 0  & 0  & 0 & \textbf{15} \\
    Eve     & 0 & 0 & 0 & 1  & 0  & 7  & 0 & \textbf{8}  \\
    \bottomrule
  \end{tabular}
\end{table}

En cada sesión con actividad solo proponen uno o dos agentes, y quiénes son cambia de una sesión a
otra. Ningún agente destaca de forma sistemática, algo esperable porque los cinco ejecutan el mismo
código: el número de propuestas depende de la instancia (activos y requisitos de cada agente), no
del agente.

\subsubsection{Acuerdos confirmados y efecto en la utilidad}
\label{sec:ej1-acuerdos}

La tabla~\ref{tab:ej1-acuerdos} detalla los 5 acuerdos confirmados: quién lo propuso (el prefijo
del ID de la propuesta), qué agentes participaron, qué intercambios incluía y cómo cambió la utilidad
(dólares producidos) de cada participante entre el inicio y el final de la sesión.

\begin{table}[htbp]
  \centering
  \small
  \caption{Acuerdos confirmados y variación de utilidad de los participantes.}
  \label{tab:ej1-acuerdos}
  \begin{tabularx}{\textwidth}{c l l Y Y}
    \toprule
    Sesión & Acuerdo & Participantes & Contenido & Ganancia de utilidad \\
    \midrule
    1 & \code{David1}  & David, Bob     & David da 6 Wood; Bob da 10 Gold
      & David $12\rightarrow14$ (+2), Bob $2\rightarrow4$ (+2) \\
    4 & \code{David34} & David, Charles & David da 4 Grain; Charles da 150 Gold (15 acciones de 10)
      & David $13\rightarrow30$ (+17), Charles $15\rightarrow17$ (+2) \\
    4 & \code{Bob7}    & Bob, Eve       & Bob da 7 Gold; Eve da 70 Iron (7 acciones de 10)
      & Bob $4\rightarrow18$ (+14), Eve $15\rightarrow17$ (+2) \\
    5 & \code{Alice13} & Alice, Bob     & Alice da 2 Gold; Bob da 130 Wood (13 acciones de 10)
      & Alice $25\rightarrow39$ (+14), Bob $5\rightarrow6$ (+1) \\
    6 & \code{Eve2}    & Eve, Bob       & Eve da 8 Grain; Bob da 20 Iron (2 acciones de 10)
      & Eve $20\rightarrow23$ (+3), Bob $22\rightarrow26$ (+4) \\
    \bottomrule
  \end{tabularx}
\end{table}

Los 5 acuerdos confirmados son bilaterales, y en la mayoría la ganancia se reparte de forma muy
desigual: en 3 de ellos el proponente gana entre 7 y 14 veces más que el otro participante (ver
sección~\ref{sec:ej1-proponente}). También destaca que 4 de los 5 acuerdos repiten la misma acción
de 10 unidades varias veces (ver sección~\ref{sec:ej1-repetidas}).

En las 7 sesiones \textbf{ningún agente terminó con menos utilidad de la inicial}, y en cada acuerdo
confirmado \textbf{todos los participantes ganan}. Es lo que el manual describe como resultado
correcto (sección~3.1 del manual).

% ---------------------------------------------------------------------
\subsection{Respuestas al ejercicio}
\label{sec:ej1-respuestas}

Antes de responder a las preguntas del enunciado conviene señalar que los ejercicios se resuelven en
orden, de modo que para obtener estos resultados no se ha aplicado ninguna mejora sobre los agentes.
Con este contexto, las respuestas son las siguientes:

\begin{itemize}
  \item \textbf{¿Cuántas propuestas ha hecho cada agente?} Depende mucho de la sesión: entre 0 y 24
        por agente y sesión (tabla~\ref{tab:ej1-por-agente}). En total hubo 86 propuestas en 7
        sesiones, y en 3 sesiones nadie propuso nada. En cada sesión con actividad solo proponen 1 o
        2 agentes. Esto se mejorará en los ejercicios 2 y 3 de la práctica.
  \item \textbf{¿Cuántas veces se ha aceptado una propuesta?} 17 mensajes \code{ACCEPT} en total
        (1, 0, 0, 12, 1, 3, 0).
  \item \textbf{¿Cuántas propuestas se han confirmado?} 5 acuerdos (1, 0, 0, 2, 1, 1, 0), además de
        1 propuesta declarada \code{ILLEGAL}. Hubo acuerdo en 4 de las 7 sesiones.
\end{itemize}

% ---------------------------------------------------------------------
\subsection{Interpretación}
\label{sec:ej1-interpretacion}

\subsubsection{Por qué en algunas sesiones no pasa nada}
\label{sec:ej1-sin-actividad}

En las sesiones sin propuestas, todos los \code{*\_nego.log} contienen solo decisiones \code{WAIT}:

\begin{lstlisting}
Decision: WAIT, at time 29675 Currently none of the deals found by me or proposed to me is selfish enough.
my aspiration level: 0.0116 opponent aspiration level: 0.9896
\end{lstlisting}

Al final de la sesión la aspiración propia $m^{\alpha}_{\alpha}(t)$ ha bajado casi a 0 y la del
oponente $m^{\alpha}_{\beta}(t)$ ha subido casi a 1. Aun así, ningún agente encontró un plan que le
diera utilidad normalizada positiva. Es decir, en esa instancia no existe ningún acuerdo rentable, o
la búsqueda no lo encontró en 30\,s (el espacio de acuerdos es enorme). El enunciado de la práctica
ya advierte de que puede ocurrir: \foreignquote{english}{you may need to run a couple of negotiation
sessions before you see any positive result}.

Los valores del log coinciden con las fórmulas (1) y (2) del manual para los grados de concesión por
defecto, $a_1 = 2$ y $a_2 = 4$:
\begin{equation}
  m^{\alpha}_{\alpha}(t) = 1 - \frac{e^{-a_1 \frac{t}{t_{dead}}} - 1}{e^{-a_1} - 1},
  \qquad
  m^{\alpha}_{\beta}(t) = \frac{e^{-a_2 \frac{t}{t_{dead}}} - 1}{e^{-a_2} - 1}.
\end{equation}
Por ejemplo, en $t = 6$\,s de 30 ($t/t_{dead} = 0{,}2$):
\begin{equation}
  m^{\alpha}_{\alpha} = 1 - \frac{e^{-0{,}4} - 1}{e^{-2} - 1} \approx 0{,}619,
  \qquad
  m^{\alpha}_{\beta} = \frac{e^{-0{,}8} - 1}{e^{-4} - 1} \approx 0{,}561.
\end{equation}
El log muestra 0,621 y 0,557. Los agentes deciden cada 300\,ms (el 1\,\% de la duración de la
sesión), así que salen 100 decisiones por agente.

\subsubsection{Muchas aceptaciones, pocas confirmaciones}
\label{sec:ej1-aceptaciones}

En el protocolo no estructurado un acuerdo solo es \textbf{vinculante cuando lo aceptan todos sus
participantes}. Proponer cuenta como aceptar.

\begin{itemize}
  \item \textbf{Sesión 4:} Charles aceptó 10 propuestas de Bob (\code{Bob21}, \code{Bob23}, \ldots),
        pero todas eran a tres bandas (\code{[Bob, Eve, Charles]}) y Eve nunca las aceptó, así que no
        se confirmó ninguna. Lo que sí se confirmó fueron dos acuerdos bilaterales (\code{David34} y
        \code{Bob7}).
  \item \textbf{Sesión 6, el ILLEGAL:} Charles propuso \code{Charles10}, que contiene exactamente las
        mismas acciones de \code{Eve2} (ya confirmado) más un intercambio propio con Eve. Eve tenía 14
        de grano y en \code{Eve2} ya había dado 8. \code{Charles10} le pide dar otras 8, lo que la
        dejaría con $-2$. Cuando Bob acepta, el Notario detecta la incompatibilidad con un acuerdo
        anterior y lo marca \textbf{ILLEGAL}. En sus logs, los agentes lo reflejan como
        \foreignquote{english}{Proposal Charles10 is inconsistent with earlier confirmed agreements}.
\end{itemize}

\subsubsection{El proponente se lleva la mayor parte de la ganancia}
\label{sec:ej1-proponente}

En 3 de los 5 acuerdos el proponente gana mucho más que quien acepta (+17 frente a +2, +14 frente a
+2, +14 frente a +1). Encaja con la estrategia NB\textsuperscript{3} (sección~2.4.1 del manual):

\begin{itemize}
  \item al \textbf{proponer}, el agente elige el plan \textbf{más egoísta} de entre los que superan
        la aspiración del oponente;
  \item al \textbf{aceptar}, basta con que el plan sea individualmente racional y supere la
        aspiración propia, que baja con el tiempo.
\end{itemize}

Quien encuentra primero el acuerdo fija un reparto favorable para sí, y el otro lo acepta porque
mejora su situación aunque sea poco. NB\textsuperscript{3} nunca rechaza: las propuestas malas
simplemente se ignoran.

\subsubsection{Acciones repetidas en los acuerdos}
\label{sec:ej1-repetidas}

Al confirmar un acuerdo, el log muestra una sola acción para un participante y la misma acción
repetida muchas veces para el otro. Por ejemplo, \code{Bob7} en la sesión~4:

\begin{lstlisting}
CommodityMarketNotary.verifyValidity() DEAL CONFIRMED: Agents: [Bob, Eve]
Bob supplies 7 units of GOLD to Eve
Eve supplies 10 units of IRON to Bob      (×7)
\end{lstlisting}

Es literal: el acuerdo tiene 8 acciones y 7 son idénticas. El motivo es que \code{getSplitLabels()},
al crear los nodos de cantidad, solo genera valores de 1 a 10 (\code{int maxValue = 10;}). Por tanto
una acción mueve como mucho 10 unidades, y para transferir más la búsqueda tiene que encadenar varias
acciones con el mismo proveedor, mercancía y consumidor. El límite lo pone el agente; el dominio
admitiría hasta 99 unidades por transacción (\code{CommodityAssets.MAX\_TRANSACTION\_VALUE}).

Que un lado tenga 1 acción y el otro 7 depende de lo que necesita cada uno
(tabla~\ref{tab:ej1-cuello}).

\begin{table}[htbp]
  \centering
  \small
  \caption{Cuello de botella de cada participante en el acuerdo \code{Bob7} (sesión 4).}
  \label{tab:ej1-cuello}
  \begin{tabular}{lllll}
    \toprule
        & Requisitos por dólar & Activos iniciales & Cuello de botella & Con el acuerdo \\
    \midrule
    Bob & 4 Oil, 5 Iron  & 88 Oil, 20 Iron  & Iron: $20/5 = 4$  & +70 Iron $\rightarrow 90/5 = \mathbf{18}$ \\
    Eve & 5 Gold, 7 Wood & 78 Gold, 151 Wood & Gold: $78/5 = 15$ & +7 Gold $\rightarrow 85/5 = \mathbf{17}$  \\
    \bottomrule
  \end{tabular}
\end{table}

Bob necesita mucho hierro (70 unidades = 7 acciones de 10), mientras que a Eve le bastan 7 unidades
de oro (1 acción). A continuación se detallan algunas consecuencias derivadas de esta situación:

\begin{itemize}
  \item \textbf{Ramas muy profundas.} Cada acción ocupa 4 niveles del árbol (proveedor
        $\rightarrow$ mercancía $\rightarrow$ cantidad $\rightarrow$ consumidor), así que \code{Bob7}
        necesita 32 niveles y \code{Alice13} (130 de madera, 14 acciones), 56. Los acuerdos que
        requieren grandes cantidades son difíciles de alcanzar en 30\,s.
  \item \textbf{Representaciones múltiples.} La misma transferencia total puede componerse de muchas
        formas (10+10+\ldots, 9+10+\ldots), y eso va en contra de la recomendación de la sección~4
        del manual de que cada acuerdo tenga una única representación en el árbol.
\end{itemize}

Es un margen de mejora claro del agente: por ejemplo, cantidades con saltos crecientes (1, 2, 5, 10,
20, 50\ldots) o ajustadas a lo que realmente necesita el consumidor.

\subsubsection{El ruido en los modelos de los oponentes}
\label{sec:ej1-ruido}

Cada agente conoce sus propios requisitos exactos, pero los de los demás le llegan \textbf{con
ruido}. Por eso calcula las cotas y la utilidad de los oponentes de forma aproximada. En estas 7
sesiones no causó pérdidas, pero sí lo observamos en una ejecución exploratoria previa (un agente
bajó de 3 a 0 de utilidad). Así que no está garantizado que el resultado sea individualmente
racional para todos.

% ---------------------------------------------------------------------
\subsection{Conclusiones}
\label{sec:ej1-conclusiones}

\begin{itemize}
  \item Al ser instancias aleatorias, los números concretos cambian de una sesión a otra. Sin
        embargo, el comportamiento cualitativo se repite: sesiones sin acuerdo, muchas más
        aceptaciones que confirmaciones, ventaja del proponente y acciones repetidas.
  \item El agente de ejemplo es funcional pero conservador. Con los grados de concesión por defecto
        y 30\,s de búsqueda, solo en algo más de la mitad de las sesiones llega a algún acuerdo. Los
        ejercicios 2 (mejor cota superior, más poda) y 3 (ajuste de los grados de concesión) atacan
        precisamente estas limitaciones.
\end{itemize}


% \input{ejercicio2}
% \input{ejercicio3}

\end{document}
